% paper.tex -- draft, A&A LaTeX class (aa.cls v7.0, EDP Sciences)
% Structure only: Equation of State, Field Equations, Castro, Braithwaite method.
% TODO: title, author list, affiliations, abstract, results/discussion sections.

\documentclass{aa}

\usepackage{graphicx}
\usepackage{amsmath}
\usepackage{amssymb}
\usepackage{txfonts}
\usepackage{natbib}

\begin{document}

   \title{Magnetized white dwarfs with mixed poloidal-toroidal fields}

   \subtitle{Draft -- structure and core physics only}

   \author{TBD\inst{1}
          }

   \institute{Affiliation TBD\\
              \email{tbd@example.org}
             }

   \date{Draft \today}

   \abstract
   % context
   {TBD.}
   % aims
   {TBD.}
   % methods
   {TBD.}
   % results
   {TBD.}
   % conclusions
   {}

   \keywords{white dwarfs -- stars: magnetic field -- magnetohydrodynamics (MHD) -- methods: numerical}

   \maketitle

%-------------------------------------------------------------------

\section{Introduction}

TBD.

%-------------------------------------------------------------------

\section{Equation of state}
\label{sec:eos}

The background stellar structure throughout this work is that of a
fully degenerate, zero-temperature electron gas
\citep[ztwd;][]{chandrasekhar1935}. Pressure and density are both
parametrized by the normalized Fermi momentum $x \equiv p_F/(m_e c)$,

\begin{align}
P(x) &= A \left[ x(2x^2-3)\sqrt{x^2+1} + 3\sinh^{-1}x \right], \\
\rho(x) &= B\,x^3,
\end{align}

\noindent where $A$ is a fixed combination of fundamental constants
($m_e$, $c$, $\hbar$) and $B$ depends on the composition only through
the mean molecular weight per electron $\mu_e = 1/Y_e$, with $Y_e$ the
number of electrons per baryon. This is a purely mechanical equation of
state: at fixed composition, pressure depends on density alone, with no
temperature dependence. Consequently, once the configuration is
restricted to be non-rotating and field-free, the entire stellar
structure (mass, radius, density profile) is determined by exactly two
physical parameters: the central density $\rho_c$ and $\mu_e$.

The specific enthalpy has a closed form,
\begin{equation}
H(x) = \frac{8A}{B}\left[\sqrt{1+x^2} - 1\right],
\end{equation}
which is what makes the self-consistent field (SCF) iteration of
Sect.~\ref{sec:field-eq} tractable: at each iteration the density is
recovered directly from the local enthalpy via the inverse relation
$x(H) = \sqrt{(1+HB/8A)^2 - 1}$, with $H \le 0$ mapped to $\rho = 0$
(the stellar surface).

Near the surface ($x \ll 1$) the equation of state reduces to
$\rho \propto H^{3/2}$, an effective polytropic index $n=3/2$; in the
relativistic core ($x \gg 1$) it approaches $\rho \propto H^3$,
$n=3$ -- the marginal polytropic index at which the equilibrium mass
becomes independent of central density, the origin of the
Chandrasekhar mass limit.

At sufficiently high density, inverse beta decay (electron capture)
changes the composition locally and the constant-$\mu_e$ cold equation
of state above ceases to describe a real white dwarf. We adopt the
tabulated threshold densities of \citet{boshkayev2013} for this
validity limit; all central densities discussed below are kept below
this threshold for the assumed composition.

%-------------------------------------------------------------------

\section{Field equations}
\label{sec:field-eq}

Axisymmetric magnetohydrostatic equilibria with a poloidal field are
governed by the Grad-Shafranov equation \citep{grad1958,shafranov1966}
for the poloidal flux function $u(\varpi,z) = \varpi A_\varphi$, where
$\varpi$ is the cylindrical radius and $A_\varphi$ the azimuthal
component of the vector potential. The poloidal field components follow
from $u$ as
\begin{equation}
B_\varpi = -\frac{1}{\varpi}\frac{\partial u}{\partial z}, \qquad
B_z = \frac{1}{\varpi}\frac{\partial u}{\partial \varpi},
\end{equation}
and the equilibrium (force balance combined with the Poisson equation
for self-gravity) reduces to a single elliptic equation for $u$ sourced
by the current distribution and, self-consistently, by $u$ itself
through the poloidal current function $M(u)$.

Following the SCF approach of \citet{hachisu1986}, extended to
magnetized configurations, the poloidal current function is taken as
linear in the flux, $M(u) = k_0\,u$, so that the Bernoulli integral
governing hydrostatic balance reads
\begin{equation}
H + \Phi - M(u) = C,
\label{eq:bernoulli}
\end{equation}
with $\Phi$ the gravitational potential and $C$ a constant fixed by the
boundary condition at the stellar surface. Equation~\ref{eq:bernoulli}
and the Grad-Shafranov equation are solved iteratively together with
the Poisson equation for $\Phi$, alternating between an update of the
density field (from the equation of state, Sect.~\ref{sec:eos}) and an
update of the flux function $u$, until convergence.

A toroidal field confined within the last closed poloidal field line
($u > u_c$) is added through a toroidal current function $\beta(u)$,
parametrized as $\beta(u) = \zeta\,(u-u_c)^{m+1}$ with $m\ge1$, giving
$B_\varphi = \beta(u)/\varpi$ inside the closed-field region and zero
outside. A purely toroidal, non-mixed branch is also considered
separately, with $B_\varphi = K\rho^m\varpi^{2m-1}$ derived directly
from the Lorentz force (Maxwell stress tensor) rather than from a flux
function; this branch requires no Grad-Shafranov solve, at the cost of
being unable to represent a poloidal component. A genuinely
self-consistent \emph{mixed} poloidal-toroidal equilibrium -- both
components simultaneously present and mutually consistent within the
same SCF solve -- is outside the scope of the axisymmetric formulation
used here; Sect.~\ref{sec:braithwaite} describes how mixed
configurations are instead obtained dynamically.

Purely toroidal and purely poloidal equilibria are each subject to a
non-axisymmetric instability of azimuthal wavenumber $m=1$
\citep{tayler1973,markey1973}, which the axisymmetric formulation of
this section cannot, by construction, test for.

%-------------------------------------------------------------------

\section{Castro}
\label{sec:castro}

Dynamical (as opposed to SCF-equilibrium) evolution is carried out with
Castro \citep{castro}, a massively parallel, block-structured
adaptive-mesh-refinement (AMR) code for compressible astrophysical
flows, built on the AMReX framework \citep{AMReX}. The hydrodynamics
solver uses an unsplit corner-transport-upwind scheme \citep{ctu} with
piecewise-parabolic reconstruction \citep{ppm}; self-gravity is
included for isolated (non-periodic) configurations, and a
constrained-transport magnetohydrodynamics (MHD) solver extends the
same integrator to evolve a magnetic field while preserving
$\nabla\!\cdot\!\mathbf{B}=0$ to machine precision by construction.
Castro accepts an arbitrary equation of state and reaction network
through the StarKiller Microphysics interface; for this work the
network is reduced to a single inert species reproducing the assumed
$\mu_e$, and the equation of state is the same zero-temperature
degenerate electron gas of Sect.~\ref{sec:eos}, so that a star built in
hydrostatic equilibrium by the SCF solver is evolved under the
identical thermodynamic relation it was constructed with.

Self-gravity for an isolated spherical mass distribution is obtained
from Castro's monopole gravity solver, which integrates the enclosed
mass along radial shells. Initial conditions are supplied to Castro as
a one-dimensional radial density (and composition) profile, sampled
onto the three-dimensional Cartesian AMR grid; the sampling itself
(cell-volume averaging versus point sampling, and the alignment of the
grid relative to the stellar center) has a measurable effect on how
closely the discretized initial condition reproduces the target central
density, which we account for when initializing a background star for
dynamical evolution.

%-------------------------------------------------------------------

\section{The Braithwaite method}
\label{sec:braithwaite}

As noted in Sect.~\ref{sec:field-eq}, the axisymmetric SCF formulation
cannot produce a genuinely self-consistent mixed poloidal-toroidal
equilibrium, nor can it test the stability of any equilibrium against
non-axisymmetric perturbations -- a serious limitation given that both
purely poloidal and purely toroidal fields are themselves
non-axisymmetrically unstable \citep{tayler1973,markey1973}. The
resolution adopted here follows \citet{braithwaite2004} and
\citet{braithwaite2006}: rather than constructing a mixed field and
separately testing its stability, a \emph{random}, multi-scale magnetic
field is seeded within an otherwise field-free, non-rotating background
star (Sect.~\ref{sec:eos}) and allowed to relax dynamically under full
three-dimensional, ideal MHD in Castro (Sect.~\ref{sec:castro}). Because
the evolution is a genuine dynamical calculation and not an imposed
equilibrium, any configuration that survives the relaxation is, by
construction, stable to the perturbations sampled by the initial random
field -- including non-axisymmetric ones the Grad-Shafranov formulation
cannot represent at all.

The random seed field is constructed as a vector potential $\mathbf{A}$
-- a superposition of Fourier-like modes spanning a band of length
scales, with random amplitude, phase, and direction for each mode --
multiplied by a smooth radial envelope that confines the field within
the stellar volume before the magnetic field itself is formed as the
discrete curl, $\mathbf{B} = \nabla\times\mathbf{A}$. Windowing the
potential rather than the field is required for
$\nabla\!\cdot\!\mathbf{B}=0$ to be preserved exactly: multiplying an
already-divergence-free $\mathbf{B}$ by a confinement window after the
curl has been taken would reintroduce a spurious divergence, whereas
curling an already-windowed potential does not. The total seed field
energy is calibrated, prior to relaxation, to a chosen fraction of the
star's gravitational binding energy, $E_{\mathrm{mag}}/|W|$.

Once released, the field is left free to evolve under the full MHD
equations; no particular final poloidal-to-toroidal ratio is imposed.
The relaxed configuration's toroidal fraction,
$E_{\mathrm{tor}}/E_{\mathrm{mag}}$, is measured directly from the
evolved field, and its dependence on the initial random seed
(reproducibility across independent realizations of the random field,
at fixed $E_{\mathrm{mag}}/|W|$) is examined empirically rather than
assumed.

\begin{acknowledgements}
TBD.
\end{acknowledgements}

%-------------------------------------------------------------------

\bibliographystyle{aa}
\bibliography{paper}

\end{document}
